\documentclass[aspectratio=169]{beamer}

\usetheme{Madrid}
\usecolortheme{default}

\usepackage[utf8]{inputenc}
\usepackage[T1]{fontenc}
\usepackage{graphicx}
\usepackage{booktabs}
\usepackage{amsmath}
\usepackage{siunitx}

\title{Mini Report: Unsupervised Refinement on Qwen3 Audio Embeddings}
\author{NXSCough Experiments}
\date{\today}

\begin{document}

\begin{frame}
  \titlepage
\end{frame}

\begin{frame}{Pipeline Summary}
\begin{enumerate}
  \item Extract embeddings with \texttt{Qwen3\_Extractor} (mean+std pooled representation).
  \item Baseline supervised evaluation across 3 classifiers.
  \item Harmonize embedding space with empirical Bayes ComBat.
  \item Detect/remove outliers (Isolation Forest + per-database centroid distance ensemble, 3\%).
  \item Estimate acoustic TB likelihood (outlier-factor score + Bayesian GMM posterior ensemble).
  \item Relabel TB samples into:
  \begin{itemize}
    \item \textbf{Acoustic-TB} (high score),
    \item \textbf{Label-only-TB} (low score), then filter label-only TB.
  \end{itemize}
\end{enumerate}
\end{frame}

\begin{frame}{Qwen3 Extractor Architecture}
\begin{itemize}
  \item Backbone: audio tower of \texttt{Qwen3-Omni-30B-A3B-Thinking}, used as a \textbf{frozen feature extractor}.
  \item Audio pipeline:
  \[
    \text{wav} \rightarrow \text{Mel frontend} \rightarrow \text{Qwen3 audio encoder (3$\times$ CNN stride)} \rightarrow \text{frame embeddings}
  \]
  \item The encoder outputs frame-level hidden states with base dimension \textbf{2048}.
  \item Temporal pooling: compute mean $\mu$ and standard deviation $\sigma$ over the time axis.
  \item Final utterance embedding:
  \[
    z = [\mu\,;\,\sigma] \in \mathbb{R}^{4096}
  \]
  Concatenating mean and std doubles the dimension to \textbf{4096}, which is used for all downstream experiments.
\end{itemize}
\begin{block}{Goal}
  Extract a fixed-length, TB-relevant audio representation from raw waveforms without any fine-tuning.
\end{block}
\end{frame}

\begin{frame}{Classifier Evaluation Protocol}
\begin{itemize}
  \item Three classifiers are compared on the same embedding space:
  \begin{itemize}
    \item Logistic Regression, MLP, SVM (RBF).
  \end{itemize}
  \item Each classifier is wrapped in a pipeline:
  \[
    \text{StandardScaler} \rightarrow \text{PCA}(50) \rightarrow \text{Classifier}
  \]
  \item Cross-validation uses \texttt{StratifiedGroupKFold} with participant ID as group key,
        ensuring no patient appears in both train and validation folds simultaneously.
  \item Training is done on CODA (DB0); external evaluation is done on CIRDZ (DB1) and \textbf{TBScreen} (DB2).
  \item Final reported score: mean AUROC across folds, then averaged across all classifiers.
\end{itemize}
\begin{block}{Goal}
  Evaluate TB detection quality at each pipeline stage using a consistent, leakage-free protocol across multiple classifiers.
\end{block}
\end{frame}

\begin{frame}{ComBat Harmonization}
\begin{itemize}
  \item ComBat is an empirical Bayes batch-effect correction method originally developed for genomics.
  \item Each database is treated as one batch. It estimates and removes:
  \begin{itemize}
    \item additive batch shift $\hat{\gamma}$ (location),
    \item multiplicative batch scale $\hat{\delta}$ (variance).
  \end{itemize}
  \item Harmonization quality is evaluated with Ridge-based predictability of database labels in PCA space:
  \begin{itemize}
    \item Raw: $R^2_{\text{db}} = 0.0198$, $R^2_{\text{dis}} = 0.0023$
    \item ComBat: $R^2_{\text{db}} \approx $\num{2.16e-8}, $R^2_{\text{dis}} = 0.0017$
  \end{itemize}
  \item After ComBat, database identity (CODA vs.\ CIRDZ) is no longer predictable from the embedding space, while disease information is partially preserved.
\end{itemize}
\begin{block}{Purpose}
  Reduce dataset-specific (batch) variation in embeddings so that downstream classifiers focus on TB-related acoustic structure rather than database identity.
\end{block}
\end{frame}

\begin{frame}{ComBat Result}
  \centering
  \includegraphics[width=0.94\linewidth,keepaspectratio]{figures/combat.png}

  \vspace{0.3em}
  \scriptsize
  \textbf{Caption:} PCA projections of embeddings before (left) and after (right) ComBat harmonization. Database-related variation is reduced after correction. External AUROC improves on CIRDZ from $0.50\pm0.05$ to $0.55\pm0.02$ and on TBScreen from $0.47\pm0.04$ to $0.48\pm0.02$.
\end{frame}

\begin{frame}{Outlier Detection Strategy}
\begin{itemize}
  \item Assumed contamination rate: \textbf{3\%}.
  \item Two complementary scores are combined:
  \begin{itemize}
    \item \textbf{Isolation Forest}: detects globally anomalous points in PCA-reduced (50-dim) embedding space.
    \item \textbf{Per-database centroid distance}: flags samples whose distance from the database centroid exceeds $\sigma_{\text{threshold}} = 2.5$ standard deviations, normalized by Median Absolute Deviation.
  \end{itemize}
  \item Weighted ensemble outlier score:
  \[
    s = 0.50\,s_{\text{IF}} + 0.50\,s_{\text{centroid}}
  \]
  \item Samples in the top 3\% of $s$ are removed before the next stage.
\end{itemize}
\begin{block}{Purpose}
  Remove samples whose embeddings are unlikely to represent clean cough signals, such as noisy recordings or non-cough acoustic content.
\end{block}
\end{frame}

\begin{frame}{Outlier Removal Result}
  \centering
  \includegraphics[width=0.94\linewidth,keepaspectratio]{figures/outlier.png}

  \vspace{0.3em}
  \scriptsize
  \textbf{Caption:} Left: detected outliers (\texttimes{}) overlaid on PCA projection. Right: retained inlier cloud after cleaning.\quad
  Removed the top \textbf{3\%} by ensemble outlier score. Mean AUROC across classifiers on CODA remains stable ($0.66\pm0.03$) while external AUROC improves on TBScreen to $0.53\pm0.02$.
\end{frame}

\begin{frame}{Acoustic TB Likelihood Modeling}
\begin{itemize}
  \item \textbf{Outlier-factor margin score (LOF)} --- a density-based ``TB-likeness'' score:
  \begin{itemize}
    \item Fit Local Outlier Factor models on TB and Non-TB subsets in PCA space.
    \item Compute a posterior-like score from the relative class densities (high score $\Rightarrow$ more TB-like).
  \end{itemize}
  \item \textbf{Bayesian GMM posterior (BGMM)} --- class-conditional density models with a Dirichlet-process prior:
  \begin{itemize}
    \item Fit separate Bayesian Gaussian mixture models on TB and Non-TB embeddings.
    \item Score $= P(\mathrm{TB}\mid x)$ via a numerically-stable posterior computed from log-likelihoods.
  \end{itemize}
  \item Final acoustic TB likelihood: simple average of LOF-margin and BGMM scores.
\end{itemize}
\begin{block}{Goal}
  Assign each sample an unsupervised score reflecting how acoustically TB-like it is in the embedding space, to enable label refinement without additional annotations.
\end{block}
\end{frame}

\begin{frame}{Acoustic TB Likelihood Result}
  \centering
  \includegraphics[width=0.90\linewidth,height=0.72\textheight,keepaspectratio]{figures/TB Likelihood.png}

  \vspace{0.25em}
  \scriptsize
  \textbf{Caption:} Top: score distributions for TB and Non-TB. Bottom: box plots with Mann--Whitney U significance test.\quad
  Scores are produced by the LOF+BGMM ensemble; TB scores are shifted higher than Non-TB scores, indicating separable acoustic structure.
\end{frame}

\begin{frame}{Relabeling Strategy}
\begin{itemize}
  \item Only original TB-labeled samples are relabeled; Non-TB samples are unchanged.
  \item Threshold $\tau$ is estimated from the \emph{1D acoustic-score distribution} using a Gaussian mixture model:
  \begin{itemize}
    \item Fit a GMM to all acoustic scores and select the best model by BIC.
    \item Choose the highest-mean component as the \textbf{Acoustic-TB} mode.
    \item Set $\tau$ at the posterior crossover point where $P(\text{Acoustic-TB}\mid s)=0.5$.
    \item Fallback: if the score distribution is unimodal, use the 95th percentile of Non-TB scores.
  \end{itemize}
  \item Interpretation by position relative to $\tau$:
  \begin{itemize}
    \item score $\geq \tau$ $\Rightarrow$ \textbf{Acoustic-TB}: embedding strongly matches the TB acoustic profile.
    \item score $< \tau$ $\Rightarrow$ \textbf{Label-only-TB}: annotated as TB but acoustically ambiguous; excluded from training.
  \end{itemize}
  \item Optional refinement: participant-level majority voting to reduce segment-level noise.
  \item This is unsupervised label refinement, not class balancing.
\end{itemize}
\begin{block}{Purpose}
  Refine noisy TB annotations by retaining only samples that are both annotated as TB \emph{and} acoustically consistent with the TB embedding cluster.
\end{block}
\end{frame}

\begin{frame}{Relabeling Result}
  \centering
  \includegraphics[width=0.90\linewidth,keepaspectratio]{figures/TB Likelihood - After plit.png}

  \vspace{0.3em}
  \scriptsize
  \textbf{Caption:} Score distribution after relabeling. Left of the dashed threshold $\tau$: Label-only-TB (acoustically ambiguous, excluded). Right of $\tau$: Acoustic-TB (kept).\quad
  The split is summarized per database (CODA/CIRDZ/TBScreen) to verify that the refinement is not dominated by a single source.
\end{frame}

\begin{frame}{Classification Performance (AUROC)}
\scriptsize
\begin{table}
\centering
\setlength{\tabcolsep}{3pt}
\renewcommand{\arraystretch}{1.15}
\resizebox{0.98\textwidth}{!}{%
\begin{tabular}{lccc}
\toprule
Stage & CODA (CV) & CIRDZ (Ext.) & TBScreen (Ext.) \\
\midrule
Baseline & $0.66\pm0.03$ & $0.50\pm0.05$ & $0.47\pm0.04$ \\
After ComBat & $0.67\pm0.03$ & $0.55\pm0.02$ & $0.48\pm0.02$ \\
After Outlier Removal & $0.66\pm0.03$ & $0.55\pm0.02$ & $0.53\pm0.02$ \\
After Acoustic-TB Filtering & $0.87\pm0.01$ & $0.71\pm0.01$ & $0.88\pm0.02$ \\
\bottomrule
\end{tabular}
}
\end{table}

\vspace{0.4em}
\footnotesize
\parbox{0.98\textwidth}{%
\textbf{Summary:} ComBat and outlier removal provide incremental gains.\\
Acoustic-TB filtering is the dominant step, improving external generalization on both CIRDZ (DB1) and TBScreen (DB2).%
}
\end{frame}

\end{document}
